\scanpage{62}
(au dessus \unsure{au quel|duquel} $X^*$ est étale \deleted{galoi} principal de groupe $\underline V$), tout comme $E$ sur $E_0$) mais non au voisinage des pts de $D_J\subset\Sigma_J$ (vu que $E_0/\Sigma_J$ y est ramifié, mais non $E/X$ en les points de $\Delta=X\mid D_J$, or si c'était cartésien il en résulterait que $E_0/\Sigma_J$ est étale \unsure{sur|en|aux} les pts de $D_J$ aussi, par descente plate\ldots). Quand $=$ il doit être facile de vérifier que le diagramme total se reconstitue \unsure{à|:} partir de $E_0$ ($\longmapsto$ \inserted{$\underline J$ et} $\Sigma=\Sigma_{\underline J}$), et $E_0\longrightarrow\Sigma$, \inserted{et section $t$ de $\Sigma^*$}) et du revêtement \deleted{principal} $X$ de $\Sigma$ avec groupe \inserted{\unsure{de type|d'op.|d'op|étales}} $\underline V$, : la donnée de $(E,\sigma)$ i.e de $(E_0,I)$ \deleted{doit être} \inserted{est} équivalente aussi à $E_0$ et du revêtement $X$ (d'où en effet \deleted{$\underline I$}$\underline K$ et $\underline I\simeq\underline K\wedge(X\mid t)^{\unsure{\prime|,}}$\ldots)

Cela montre que finalement $(E,\sigma)$ i.e. le $I$ est un peu de la poudre aux yeux, on va prendre par la suite $E=E_0$, $\sigma=\sigma_0$, donc $I=\underline V={}_2E_0=M_0$, torseur trivial. On a donc \deleted{$\unsure{\Sigma|J}$} $X=\Sigma_J$ et on trouve donc
% archive sheet number: 61
